\documentclass{article}
\begin{document}
This note records three local linear-algebra identities on a two-index
polynomial and a $2\times 2$ coordinate projector. It is a synthetic
public demo, not a research manuscript.

Define the two-index polynomial
\begin{equation}
\label{eq:K-def}
K(m,n) := m+n.
\end{equation}
Swapping the two arguments is a dummy relabeling
\begin{equation}
\label{eq:K-swap}
K(n,m) = n+m,
\end{equation}
and therefore $K(m,n)=K(n,m)$.

Let $P$ be the $2\times 2$ matrix that keeps the second coordinate and
discards the first,
\begin{equation}
\label{eq:projector}
P = \begin{pmatrix} 0 & 0 \\ 0 & 1 \end{pmatrix},
\qquad
P^{2}-P = 0.
\end{equation}

For a generic two-index array $A(m,n)$, the local pair equals twice the
symmetric part $S=(A(m,n)+A(n,m))/2$,
\begin{equation}
\label{eq:pairwise}
A(m,n)+A(n,m)-2S = 0.
\end{equation}

The three local identities are collected on one worksheet; that assembly
step is bookkeeping, not an additional exact residual.
\end{document}
